\documentclass[10pt,a4paper]{amsart}
\usepackage[margin=19mm]{geometry}
\usepackage{amsmath,amssymb,mathtools}
\usepackage[scr=rsfs,cal=euler]{mathalfa}
\usepackage{xfrac}
\usepackage[unicode,hidelinks,bookmarks=false]{hyperref}
\newcommand{\ie}{i.e.\ }
\newcommand{\cf}{cf.\ }
\newcommand{\resp}{respectively\ }
\newcommand{\Subsection}{\subsection}
\providecommand{\nobreakdash}{\nobreak\hbox{-}}
\setlength{\emergencystretch}{2em}
\multlinegap0pt
\usepackage{amssymb}
\usepackage{tikz}
\usepackage{multirow}
\usepackage{booktabs}
\newtheorem{theorem}{Theorem}
\title{Node-Kayles on Trees}
\author{Nuttanon Songsuwan}
\date{Source: arXiv:2512.24221v1 (CC BY 4.0)}
\begin{document}
\maketitle

\noindent\emph{Source and licence.} Node-Kayles on Trees. Version arXiv:2512.24221v1 (CC BY 4.0), distributed by arXiv under the Creative Commons Attribution 4.0 International licence, http://creativecommons.org/licenses/by/4.0/, as recorded in the arXiv licence metadata for this exact version and shown on the abstract page https://arxiv.org/abs/2512.24221v1. Full licence notice: \texttt{licenses/arxiv-2512.24221-CC-BY-4.0.txt}.\par\smallskip

\noindent\textit{Editorial scope.} Counted unit: the article's theorem Tnm (source0.tex lines 278--288). The article's complete original proof of it is delivered with it (source0.tex lines 289--356, one \texttt{proof} environment of 68 lines). No mathematics is written, repaired or re-derived: the statement and the proof are the article's own lines, in the article's own notation, and no retained line is substituted or re-flowed.  No further source-local context is needed: every cross-reference of every retained line resolves inside the two delivered spans.  Nothing was omitted from any retained span, so the package carries no omission record.  No retained line contains a citation, so the wrapper carries no bibliography and no bibliography entry.  No retained line contains a figure environment, an image-inclusion command or drawing code, so no image is delivered and the package declares no asset dependency.  Editorial formatting only, in addition: an AMS \texttt{amsart} preamble that restates the article's own declarations and macros used by the retained lines, namely the environments \texttt{theorem} on separate counters, together with the standard packages those lines need.  The retained environments print in this document as Theorem 1 in order of appearance, whereas the article prints the same items with the numbering of its own full document.  The complete native source of the article as distributed in the arXiv e-print for this exact version is bundled unchanged as \texttt{sources/191865/source0.tex}.
\begin{theorem}\label{Tnm}
	The Grundy value sequence $\left( \mathcal{G}\left(T_{\{n,m\}}\right) : n,m \in \mathbb{N} \right)$ are given by the following:
	\begin{equation*}
		\mathcal{G}\left(T_{\{n,m\}}\right) = \begin{cases}
			1,	\quad \text{if $m$ is even};\\
			2,	\quad \text{if $n,m$ are odd};\\
			3,	\quad \text{if $n$ is even and $m$ is odd}
		\end{cases}
	\end{equation*}
	for all $n \in \mathbb{N}$
\end{theorem}

\begin{proof}
	The move of the first player is consist of 3 action
	\begin{itemize}
		\item play at the root: $\bigoplus_{i=1}^{n\cdot m}\mathcal{G}\left(\cdot\right)$
		\item play at the first level: $\bigoplus_{i=1}^{n-1}\mathcal{G}\left(T_{\{m\}}\right)$
		\item play at the second level: $\mathcal{G}\left(T_{\{n-1,m\}}\right) \oplus \bigoplus_{i=1}^{m-1}\mathcal{G}\left(\cdot\right)$
	\end{itemize}
	By the definition,
	\begin{equation*}
		\mathcal{G}\left(T_{\{n,m\}}\right) = \mathrm{mex}\left\{\bigoplus_{i=1}^{n\cdot m}\mathcal{G}\left(\cdot\right), \bigoplus_{i=1}^{n-1}\mathcal{G}\left(T_{\{m\}}\right), \mathcal{G}\left(T_{\{n-1,m\}}\right) \oplus \bigoplus_{i=1}^{m-1}\mathcal{G}\left(\cdot\right)\right\}
	\end{equation*}
	Suppose $m$ is even. It is easy to see that $\mathcal{G}\left(T_{\{1,m\}}\right) = \mathcal{G}\left(T_{\{m+1\}}\right) =  1$, $\bigoplus_{i=1}^{n\cdot m}\mathcal{G}\left(\cdot\right) = 0$ and
	\begin{equation*}
		\bigoplus_{i=1}^{n-1}\mathcal{G}\left(T_{\{m\}}\right) = \begin{cases}
			0,	\quad \text{if $n$ is odd};\\
			2,	\quad \text{if $n$ is even}.
		\end{cases}
	\end{equation*} By using induction,
	\begin{align*}
		\mathcal{G}\left(T_{\{n,m\}}\right) &= \mathrm{mex}\left\{\bigoplus_{i=1}^{n\cdot m}\mathcal{G}\left(\cdot\right), \bigoplus_{i=1}^{n-1}\mathcal{G}\left(T_{\{m\}}\right), \mathcal{G}\left(T_{\{n-1,m\}}\right) \oplus \bigoplus_{i=1}^{m-1}\mathcal{G}\left(\cdot\right)\right\}\\
		&= \mathrm{mex}\left\{0, (0 \text{ or } 2), \mathcal{G}\left(T_{\{n-1,m\}}\right) \oplus  1\right\}\\
		&= \mathrm{mex}\left\{0, (0 \text{ or } 2), 1 \oplus  1\right\}\\
		&= \mathrm{mex}\left\{0, (0 \text{ or } 2)\right\}\\
		&= 1.
	\end{align*}
	If $m$ is odd, then $\mathcal{G}\left(T_{\{1,m\}}\right) = \mathcal{G}\left(T_{\{m+1\}}\right) =  2$,
	\begin{equation*}
			\bigoplus_{i=1}^{n\cdot m}\mathcal{G}\left(\cdot\right) = \begin{cases}
				0,	\quad \text{if $n$ is even};\\
				1,	\quad \text{if $n$ is odd}
			\end{cases}
	\end{equation*}
	and
	\begin{equation*}
		\bigoplus_{i=1}^{n-1}\mathcal{G}\left(T_{\{m\}}\right) = \begin{cases}
			0,	\quad \text{if $n$ is odd};\\
			1,	\quad \text{if $n$ is even}.
		\end{cases}
	\end{equation*}
	We left to show that $\mathcal{G}\left(T_{\{2,m\}}\right) =  3$. The Grundy value of $\mathcal{G}\left(T_{\{2,m\}}\right)$ is
	\begin{align*}
		\mathcal{G}\left(T_{\{2,m\}}\right) &= \mathrm{mex}\left\{\bigoplus_{i=1}^{2 m}\mathcal{G}\left(\cdot\right), \mathcal{G}\left(T_{\{m\}}\right), \mathcal{G}\left(T_{\{1,m\}}\right) \oplus \bigoplus_{i=1}^{m-1}\mathcal{G}\left(\cdot\right)\right\}\\
		&= \mathrm{mex}\left\{0, 1, 2 \oplus 0\right\}\\
		&= \mathrm{mex}\left\{0, 1, 2 \right\}\\
		&= 3
	\end{align*}
	We are ready to finish our proof by strong induction. Assume that for some integer $n \ge 2$, 
	\begin{equation*}
		\mathcal{G}\left(T_{\{i,m\}}\right) = \begin{cases}
			2,	\quad \text{if $i$ are odd};\\
			3,	\quad \text{if $i$ is even}
		\end{cases}
	\end{equation*}
	for all integers $1 \le i < n$. If $n$ is odd, we get
	\begin{align*}
		\mathcal{G}\left(T_{\{n,m\}}\right) &= \mathrm{mex}\left\{\bigoplus_{i=1}^{n\cdot m}\mathcal{G}\left(\cdot\right), \bigoplus_{i=1}^{n-1}\mathcal{G}\left(T_{\{m\}}\right), \mathcal{G}\left(T_{\{n-1,m\}}\right) \oplus \bigoplus_{i=1}^{m-1}\mathcal{G}\left(\cdot\right)\right\}\\
		&= \mathrm{mex}\left\{1,0,3 \oplus 0\right\}\\
		&= \mathrm{mex}\left\{0,1,3\right\}\\
		&= 2
	\end{align*}
	In the similar way for $n$ is even,
	\begin{align*}
		\mathcal{G}\left(T_{\{n,m\}}\right) &= \mathrm{mex}\left\{\bigoplus_{i=1}^{n\cdot m}\mathcal{G}\left(\cdot\right), \bigoplus_{i=1}^{n-1}\mathcal{G}\left(T_{\{m\}}\right), \mathcal{G}\left(T_{\{n-1,m\}}\right) \oplus \bigoplus_{i=1}^{m-1}\mathcal{G}\left(\cdot\right)\right\}\\
		&= \mathrm{mex}\left\{0,1,2 \oplus 0\right\}\\
		&= \mathrm{mex}\left\{0,1,2\right\}\\
		&= 3
	\end{align*}
\end{proof}

\end{document}
